\section{一つの量$h$の証明と性質}
\label{sec:app_h19}
記号は\secref{h19def}のとおりで，
\begin{equation*}
M=(D+\gamma I-\lambda W)_{FF},\qquad\rho_x=\frac{\lambda\,\delta_x}{\delta_x+\gamma}
\end{equation*}
とする．

\textbf{\TFaces{}の証明．}

\textit{(i)から(ii)．}$x\in F$で最初の滞在時間$T_1\sim\mathrm{Exp}(\delta_x)$を生き延びる確率は
\begin{equation*}
\mathbb{E}\,e^{-\gamma T_1}=\int_0^{\infty}\delta_x e^{-\delta_x t}\,e^{-\gamma t}\,dt=\frac{\delta_x}{\delta_x+\gamma}
\end{equation*}
であり，続くジャンプを生き延びて$j$へ移る確率は$\lambda P_{xj}$である．よって
\begin{equation*}
h(x)=\frac{\delta_x}{\delta_x+\gamma}\sum_j\lambda P_{xj}\,h(j)=\rho_x\,(Ph)(x)=\frac{\lambda\,(Wh)(x)}{\delta_x+\gamma}\qquad(x\in F)
\end{equation*}
であり，境界の値（$\mathcal{S}$で1，$\hat{Q}$で0）を代入すれば
\begin{equation*}
(\delta_x+\gamma)\,h(x)-\lambda\sum_{j\in F}W_{xj}\,h(j)=\lambda\sum_{j\in\mathcal{S}}W_{xj}\qquad(x\in F),
\end{equation*}
すなわち(ii)の$Mh_F=\lambda W_{F\mathcal{S}}\mathbf{1}$になる．

\textit{一意性．}$M$の非対角成分の行和は
\begin{equation*}
\lambda\sum_{j\in F}W_{xj}\le\lambda\,\delta_x
\end{equation*}
なので，$\gamma>0$，または$\lambda<1$かつ$\delta_x>0$なら$M$は狭義対角優位であり，対称なので正定値である（$\lambda{=}1,\gamma{=}0$では，$F$と交わる各連結成分が境界を含めば既約対角優位）．よって解は一意で，確率\eqnref{h}（方程式の最小の非負解）と一致する．

\textit{(iii)．}$\mathcal{E}$を$f_x$（$x\in F$）で微分すると
\begin{equation*}
\frac{\partial\mathcal{E}}{\partial f_x}=2\Big[(\delta_x+\gamma)f_x-\lambda\sum_jW_{xj}f_j\Big]
\end{equation*}
で，停留条件は(ii)と同じ式であり，ヘッセ行列$2M$は正定値なのでその解が最小点である．回路ではこれが自由な節点でのキルヒホッフの電流則
\begin{equation*}
\sum_j\lambda W_{xj}\,(f_x-f_j)+\big((1-\lambda)\delta_x+\gamma\big)f_x=0
\end{equation*}
にあたる．$\square$

\appblock
\textbf{三つの行き先．}$\mathcal{S}$と$\hat{Q}$の役割を入れ替えた量を$\check h$，どちらにも着く前に消える確率を$d$とすると，$M$が正則なとき
\begin{equation*}
h+\check h+d=1
\end{equation*}
である（$\lambda{=}1,\gamma{=}0$では$\check h=1-h$）．

\textit{証明．}最初の滞在と最初のジャンプで場合を分けると
\begin{equation*}
d(x)=\frac{\gamma}{\delta_x+\gamma}+\frac{\delta_x}{\delta_x+\gamma}\Big[(1-\lambda)+\lambda\sum_{j\in F}P_{xj}\,d(j)\Big]\qquad(x\in F),
\end{equation*}
すなわち$Md_F=\gamma\mathbf{1}+(1-\lambda)\boldsymbol{\delta}_F$である．一方
\begin{align*}
M\mathbf{1}_F&=(1-\lambda)\boldsymbol{\delta}_F+\gamma\mathbf{1}+\lambda W_{F\mathcal{S}}\mathbf{1}+\lambda W_{F\hat{Q}}\mathbf{1},\\
Mh_F&=\lambda W_{F\mathcal{S}}\mathbf{1},\qquad M\check h_F=\lambda W_{F\hat{Q}}\mathbf{1}
\end{align*}
だから
\begin{equation*}
M\big(h+\check h+d\big)_F=M\mathbf{1}_F
\end{equation*}
であり，$M$が正則なので$h+\check h+d=1$である．$\square$

\appblock
\textbf{\TLimits{}の証明．}

\textit{(i)．}種がないので$h(x)=\mathbb{E}_x[\lambda^{N}e^{-\gamma T}]$（$N=N_{\mathcal{S}}$，$T=T_{\mathcal{S}}$）であり，
\begin{equation*}
\lambda^{N}e^{-\gamma T}=\exp\big(N\log(1-\mu\varepsilon)-\nu\varepsilon T\big)
\end{equation*}
である．$1-e^{-u}\le u$と$-\log(1-v)\le2v$（$0\le v\le1/2$）から，$\mu\varepsilon\le1/2$のとき
\begin{equation*}
0\le\frac{1-\lambda^{N}e^{-\gamma T}}{\varepsilon}\le2\mu N+\nu T
\end{equation*}
である．有限のグラフで各連結成分が支持画像を含めば$N$と$T$は指数的な裾を持ち可積分なので，各点収束
\begin{equation*}
\frac{1-\lambda^{N}e^{-\gamma T}}{\varepsilon}\;\longrightarrow\;\mu N+\nu T\qquad(\varepsilon\to0)
\end{equation*}
に優収束定理を使えば
\begin{equation*}
\frac{1-h(x)}{\varepsilon}\;\longrightarrow\;\mu\,\mathbb{E}_x[N_{\mathcal{S}}]+\nu\,\mathbb{E}_x[T_{\mathcal{S}}]
\end{equation*}
を得る．期待値$n=\mathbb{E}[N_{\mathcal{S}}]$と$t=\mathbb{E}[T_{\mathcal{S}}]$は，最初の1歩で分けた方程式
\begin{equation*}
(D-W)_{FF}\,n=\boldsymbol{\delta}_F,\qquad(D-W)_{FF}\,t=\mathbf{1}
\end{equation*}
の解である．

\textit{(ii)．}$\hat W=D^{-1/2}WD^{-1/2}$を固有分解して
\begin{equation*}
\hat W=\sum_k\theta_k\,\phi_k\phi_k^{\top},\qquad\theta_1=1,\qquad\phi_1=\frac{D^{1/2}\mathbf{1}}{\|D^{1/2}\mathbf{1}\|}
\end{equation*}
とする（連結なら$\theta_1$は単純）．
\begin{equation*}
(1-\lambda)(I-\lambda\hat W)^{-1}=\sum_k\frac{1-\lambda}{1-\lambda\theta_k}\,\phi_k\phi_k^{\top}
\end{equation*}
の係数は，$\lambda\to1$で$k\ge2$では0に（$\theta_k=-1$でも$(1-\lambda)/(1+\lambda)\to0$），$k=1$では1に収束する．よって
\begin{equation*}
u\;\longrightarrow\;\phi_1\phi_1^{\top}\mathbf{1}_{\mathcal{S}}=\sqrt{\boldsymbol{\delta}}\;\frac{\sum_{j\in\mathcal{S}}\sqrt{\delta_j}}{\sum_i\delta_i}
\end{equation*}
である．大きなランダム近傍グラフでは1点への到達時間が次数で決まる量に退化することが知られ~\cite{vonluxburg2014hitting}，$\lambda{=}1,\gamma{=}0$の調和関数もラベルなしの点が非常に多い極限で平坦になり得る~\cite{nadler2009limit}．本稿が極限ではなく有限の代価を使うのはこのためで，(i)は極限での読み方である．$\square$

\appblock
\textbf{\TValid{}の証明．}
\begin{itemize}
\item[(a)] \textit{置換不変性．}$s$は，ラベル付きの$\mathcal{S}$，$\mathcal{R}$と，$\mathcal{U}$の画像の特徴の集まり（順序なし）だけの関数で，$\mathcal{U}$の中でどの画像を較正画像と呼ぶかを入れ替えても各画像の値は変わらない．グラフは節点の名札に依らず，$h^{(0)}$は$\mathcal{S}$を吸収壁とする同じ方程式の解，$p_{\mathcal{R}}$は固定した$\mathcal{R}$との比較，BH法はプールの$p_{\mathcal{R}}$の多重集合から閾値を決めるので種の選び方は$\mathcal{U}$の画像を対称に扱い，第2段の$h$と$s$も同様だからである．
\item[(b)] \textit{(i)と(ii)．}(a)はAdaDetect~\cite{marandon2024adadetect}の条件(8)であり，$p$は同論文の経験$p$値(10)で，同点を独立な一様乱数で崩せば同点なしの仮定も満たす（崩さずに$\le$で数えると(i)は保守的になる）．よって同論文の定理3.3が(i)とPRDSを，系3.5と3.7が(ii)を与える．系3.7を当てはめるStorey型の$\hat\pi_0$は上限を付けない推定量で，その閾値は
\begin{equation*}
\lambda_{\mathrm{S}}=\frac{\lfloor|\mathcal{C}|/2\rfloor}{|\mathcal{C}|+1}
\end{equation*}
である．
\item[(c)] \textit{(iii)．}$h^{(0)}$は$\mathcal{R}\cup\mathcal{U}$の中の置換に不変なので，$\mathcal{R}$を較正標本として同じ定理による．
\end{itemize}
クラスごとに位置で固定した分割は$\mathcal{C}$と$\mathcal{R}$を層化するので，交換可能性は厳密には仮定である．$\square$

\appblock
\textbf{\TEscape{}の証明．}

\textit{出口の分解．}$Q$を初めて出る時刻を$\tau$とする．$Q$は支持画像も種も含まないので$\tau\le\tau_{\mathcal{S}}$かつ$\tau\le\tau_{\hat{Q}}$であり，割引$\lambda^{N}e^{-\gamma T}$は$\tau$の前後の積に分かれる．強マルコフ性から
\begin{equation*}
h(y)=\mathbb{E}_y\big[\lambda^{N_\tau}e^{-\gamma\tau}\,h(X_\tau)\big]=\sum_{z\in\partial Q}H_Q(y,z)\,h(z)
\end{equation*}
で，種の上で$h=0$，それ以外で$h\le1$だから
\begin{equation*}
h(y)\le\sum_{z\in\partial Q\setminus\hat{Q}}H_Q(y,z)=\pi_Q(y)
\end{equation*}
である．

\textit{逃げる確率の上界．}最初の滞在と最初のジャンプで分けると，$\pi=\pi_Q$は写像
\begin{equation*}
\mathcal{K}(f)(y)=\rho_y\Big(e_y+\sum_{j\in Q}P_{yj}\,f(j)\Big)\qquad(y\in Q)
\end{equation*}
の最小の非負の不動点で，$0$から反復して得られる（$\pi=\lim_{k\to\infty}\mathcal{K}^k(0)$）．$\sum_{j\in Q}P_{yj}=1-e_y-\zeta_y$だから，定数$c\in[0,1]$に対し
\begin{equation*}
\mathcal{K}(c\mathbf{1})(y)=\rho_y\big(e_y+(1-e_y-\zeta_y)\,c\big)=\rho_y\big(c+(1-c)\,e_y-c\,\zeta_y\big)
\end{equation*}
で，これは$\rho_y$と$e_y$について増加，$\zeta_y$について減少する．$\bar\varphi=\varphi(\max_Q\rho,\max_Qe,\min_Q\zeta)$は
\begin{equation*}
\bar\varphi=\max_Q\rho\,\Big(\max_Qe+\big(1-\max_Qe-\min_Q\zeta\big)\,\bar\varphi\Big)
\end{equation*}
を満たすので$\mathcal{K}(\bar\varphi\mathbf{1})\le\bar\varphi\mathbf{1}$（優解）である．$\mathcal{K}$は単調（$f\le g$なら$\mathcal{K}(f)\le\mathcal{K}(g)$）なので，反復の各段で
\begin{equation*}
\mathcal{K}^k(0)\le\mathcal{K}^k(\bar\varphi\mathbf{1})\le\bar\varphi\mathbf{1}
\end{equation*}
が保たれ，$\pi_Q\le\bar\varphi$である．$\square$

\appblock
\textbf{\TCapture{}の証明．}$B$の中だけを通って支持画像に着く経路は種を通らないので$\kappa_B\le h$である．$\kappa_B$は
\begin{equation*}
\mathcal{K}_B(f)(y)=\rho_y\Big(\omega_y+\sum_{j\in B}P_{yj}\,f(j)\Big)\qquad(y\in B)
\end{equation*}
の不動点で，$\sum_{j\in B}P_{yj}=1-\omega_y-\eta_y$である．\TEscape{}と逆向きに，$\underline\varphi=\varphi(\min_B\rho,\min_B\omega,\max_B\eta)$は
\begin{align*}
\mathcal{K}_B(\underline\varphi\mathbf{1})(y)
&=\rho_y\big(\underline\varphi+(1-\underline\varphi)\,\omega_y-\underline\varphi\,\eta_y\big)\\
&\ge\min_B\rho\,\Big(\min_B\omega+\big(1-\min_B\omega-\max_B\eta\big)\,\underline\varphi\Big)=\underline\varphi
\end{align*}
を満たす（劣解）．よって$\mathcal{K}_B^k(\underline\varphi\mathbf{1})$は$k$について増加して不動点に収束する．$\max_B\rho<1$なら$\mathcal{K}_B$は縮小写像で不動点は$\kappa_B$ただ一つであり，$\lambda{=}1,\gamma{=}0$でも$B$の各成分が出口を持てば一意で，出口のない成分があれば$\min_B\omega=0$で主張は自明である．最後の主張は，$\mathcal{S}$で$h=1$，他で$h\ge0$なので
\begin{equation*}
s(x)=\rho_x\sum_jP_{xj}\,h(j)\ge\rho_x\sum_{j\in\mathcal{S}}P_{xj}=\rho_x\,\omega_x
\end{equation*}
による．$\square$

\appblock
\textbf{\CSep{}の証明．}$y$は種でないので$s(y)=h(y)\le\pi_Q(y)$である．$s(x')\le s(y)$なら$\ell(x')\le s(x')\le s(y)\le\pi_Q(y)$なので
\begin{equation*}
\{x'\in\mathcal{C}:s(x')\le s(y)\}\subseteq\{x'\in\mathcal{C}:\ell(x')\le\pi_Q(y)\}
\end{equation*}
であり，これを\eqnref{p}に代入すれば\eqnref{sep}を得る．$\square$

\appblock
\textbf{クラスに依らない．}IDのクラスごとの支持画像$\mathcal{S}_c$について，最初に着く支持画像が$\mathcal{S}_c$に属する確率を$h_c$とすると
\begin{equation*}
h=\sum_ch_c
\end{equation*}
である（事象$\{\tau_{\mathcal{S}}<\tau_{\hat{Q}}\}$は最初に着く支持画像のクラスで排反に分かれ，割引は共通）．$h$はクラスの分け方に依らない．同じグラフでクラスラベルを伝播させて最大値で読むLGC-maxでは，近いIDクラスが複数ある画像の到達確率がそれらに割れて最大値が薄まるが，$h$はその総和で割れても変わらない．

\appblock
\textbf{粒度の床．}$p\ge1/(|\mathcal{C}|+1)$で，$p$は$1/(|\mathcal{C}|+1)$刻みの値を取る．種が選ばれるなら，その数は少なくとも$\hat\pi_0m/(q(|\mathcal{R}|+1))$（プール$\mathcal{U}$の大きさ$m$の割合$\hat\pi_0/(q(|\mathcal{R}|+1))$）である．BH法の棄却数$\hat k$は$p_{(\hat k)}\le q\hat k/(\hat\pi_0m)$を満たし，$p_{\mathcal{R}}$はつねに$1/(|\mathcal{R}|+1)$以上なので
\begin{equation*}
\frac{1}{|\mathcal{R}|+1}\le p_{(\hat k)}\le\frac{q\,\hat k}{\hat\pi_0\,m},
\qquad\text{すなわち}\qquad
\hat k\ge\frac{\hat\pi_0\,m}{q\,(|\mathcal{R}|+1)}
\end{equation*}
である．ImageNet-1K（$|\mathcal{R}|{=}2{,}000$，$q{=}0.2$）ではプールの約0.25\%であり，共形$p$値の分解能から来る構造である．
